\section{Descent through a family of integral maps}
\label{sec:map-descent}

Here is the descent inequality used in the comparison of \(L_-\)
with \(L_d\). Write \(F\) for a totally complex number field and
\(v\) for one archimedean place from each conjugate pair. Let \(m\)
be the number of such places. Give the finite-dimensional \(F\)-spaces
\(V,W\) positive definite hermitian norms at each place. We use
the degree \(\widehat d\) in the squared-norm normalization of the
preceding section. In particular the degree of a free line generated
by \(x\) is \(-m^{-1}\sum_v\log\|x\|_v^2\). For an integral lattice
\(C\) in its \(F\)-span put
\(\mu(C)=\widehat d(C)/\operatorname{rk}_F(C)\) when its rank is positive.

Let \(A\subset V,\ B\subset W\) be finitely generated
\(\mathcal O_F\)-submodules, \(A\ne0\). Suppose \(n\) linear maps
\(\phi_j:V\to W\) send \(A\) into \(B\), have zero common kernel and
obey, for all \(v\) and \(x\in V\),
\begin{equation}
 \sum_{j=1}^n\|\phi_jx\|_v^2\ \le\ e^{2P}\|x\|_v^2 .
 \label{eq:descent-map-bound}
\end{equation}
Suppose \(B\) is spanned over \(F\) by elements of \(B\) of norm at
most \(e^\Lambda\) at every place. Put \(a=\operatorname{rk}_F A\),
\(N=n\operatorname{rk}_F B\). If \(\Lambda,\epsilon\ge0\),
\(N\le(1+\epsilon)a\) and \(t\le\mu(A)-2P\), then
\begin{equation}
                  \mu(B)\ \ge\ t-\epsilon(|t|+2\Lambda).
 \label{eq:descent-result}
\end{equation}
Neither a positive \(t\) nor a free lattice \(A\) is required.

\paragraph{The lower bound on quotient slopes.}
We give the determinant argument underlying the estimate for the
minimal slope. Let \(C\) be spanned over \(F\) by integral vectors of
norm at most \(e^\Lambda\) at all places. For a saturated sublattice
\(D\subset C\), put \(Q=C/D\) with its quotient norms. The images of
these vectors span \(Q\) over \(F\), and the quotient norms do not
increase. If \(r=\operatorname{rk}_F Q>0\), select \(r\) images
\(q_1,\ldots,q_r\) forming an \(F\)-basis. Their integral span
\(Q_0\) is free and of finite index in \(Q\). At every \(v\) Hadamard's
inequality gives
\[
  \det(\langle q_i,q_j\rangle_v)_{i,j}
        \ \le\ \prod_i\|q_i\|_v^2 \ \le\ e^{2r\Lambda}.
\]
The determinant formula for the degree gives
\(\widehat d(Q_0)\ge-2r\Lambda\). Enlarging a full-rank lattice
by finite index increases its degree. Indeed in a common real
realization the volume of the smaller lattice is the index times
the volume of the larger one, and degree uses the negative logarithm
of volume with positive normalization. Therefore
\(\widehat d(Q)\ge-2r\Lambda\). This also holds at rank zero, with
degree zero. Equivalently all nonzero quotient slopes are at least
\(-2\Lambda\). In the formal development the same conclusion is
obtained by putting the short lines into a saturated step of the
Harder--Narasimhan filtration. The full-span property gives full rank
of that step. A saturated sublattice of full rank equals the ambient
lattice.\footnote{\href{https://github.com/Deicyde/riemann-hypothesis/blob/40e56bab/src/Formalism/MeanHeightDescent/AverageDegreeBounds/SpanMuMin.lean}{Short generators, line heights and the saturated filtration step}.}

\paragraph{The image lattice.}
Give \(B^n\) the orthogonal-sum metric and put
\(\Phi x=(\phi_1x,\ldots,\phi_nx)\).
The common-kernel hypothesis gives injectivity. Choosing a nonzero
element of \(A\) shows that \(n>0\) and \(B\ne0\). Write \(I=\Phi A\),
and \(D=B^n\cap (F I)\) for its saturation in \(B^n\).
Then \(a=\operatorname{rk}_F I=\operatorname{rk}_F D\le N\), and
\(\widehat d(D)\ge\widehat d(I)\). By
\eqref{eq:descent-map-bound} the operator norm of \(\Phi\) at each
place is at most \(e^P\). The bound extends to the completion at that
place by density of \(F\) and continuity of the linear map.
On the top exterior power of \(F A\) the norm multiplier is at most
\(e^{aP}\), since each singular value on this subspace is at most
\(e^P\). The induced map on the determinant modules of \(A\) and
\(I\) is an integral isomorphism, so no index arises there. The
change of squared determinant norms contributes at most \(2aP\)
on averaging their logarithms. Hence
\[
             \widehat d(D)\ge\widehat d(I)
             \ge\widehat d(A)-2aP \ge at .
\]
The short generators of \(B\), placed one at a time in each coordinate,
give integral short generators of \(B^n\) over \(F\). The preceding
quotient bound yields \(\widehat d(B^n/D)\ge-2(N-a)\Lambda\).
In a saturated exact sequence, determinant degrees add when the
quotient has its quotient norm. This follows by wedging a basis of
the subspace with orthogonal lifts of a basis of the quotient at each
place and using the integral determinant isomorphism. In particular
it applies even when \(D\) is only projective over \(\mathcal O_F\).
Since the degree of \(B^n\) is \(n\widehat d(B)\), we obtain
\[
                 N\mu(B)\ \ge\ at-2(N-a)\Lambda .
\]
Subtract the right side from \(Nt\). The difference is at most
\((N-a)(|t|+2\Lambda)\). The rank hypothesis gives
\(N-a\le\epsilon a\le\epsilon N\), and \(N>0\).
Division by \(N\) proves \eqref{eq:descent-result}.
\footnote{\href{https://github.com/Deicyde/riemann-hypothesis/blob/40e56bab/src/Formalism/MeanHeightDescent/AverageDegreeBounds/DescentSlope.lean}{Slope inequality for an injective map};
\href{https://github.com/Deicyde/riemann-hypothesis/blob/40e56bab/src/Development/AbelianVarietyTheory/Proofs/DivisorProofs/HtMapLeOfPoints.lean}{direct-sum norm, comparison of ranks and final inequality}.}
In the preceding section the source takes \(A=L_-\), \(B=L_d\),
\(P=b,\ \Lambda=K,\ \epsilon=b/K\). Thus the inequality holds for
the actual ranks of the lattices on each window once the listed maps,
short generators and rank estimates have been supplied there.
